\section{Presheaves}
\label{section-presheaves}


\begin{definition}
\label{definition-presheaf}
Let $X$ be a topological space.
\begin{enumerate}
\item A {\it presheaf $\mathcal{F}$ of sets on $X$} is a rule which
assigns to each open $U \subset X$ a set $\mathcal{F}(U)$ and
to each inclusion $V \subset U$ a map
$\rho^U_V : \mathcal{F}(U) \to \mathcal{F}(V)$ such that
$\rho^U_U = \text{id}_{\mathcal{F}(U)}$ and
whenever $W \subset V \subset U$ we have
$\rho^U_W = \rho^V_W \circ \rho ^U_V$.
\item A {\it morphism $\varphi : \mathcal{F} \to \mathcal{G}$
of presheaves of sets on $X$} is a rule which assigns to each
open $U \subset X$ a map of sets $\varphi : \mathcal{F}(U)
\to \mathcal{G}(U)$ compatible with restriction maps,
i.e., whenever $V \subset U \subset X$ are open the
diagram
$$
\xymatrix{
\mathcal{F}(U) \ar[r]^\varphi \ar[d]^{\rho^U_V} &
\mathcal{G}(U) \ar[d]^{\rho^U_V} \\
\mathcal{F}(V) \ar[r]^\varphi & \mathcal{G}(V)
}
$$
commutes.
\item The category of presheaves of sets on $X$ will be denoted
$\textit{PSh}(X)$.
\end{enumerate}
\end{definition}

\noindent
The elements of the set $\mathcal{F}(U)$ are called
the {\it sections} of $\mathcal{F}$ over $U$.
For every $V \subset U$ the map
$\rho^U_V : \mathcal{F}(U) \to \mathcal{F}(V)$
is called the {\it restriction map}. We will use the
notation $s|_V := \rho^U_V(s)$ if $s\in \mathcal{F}(U)$.
This notation is consistent with the notion of restriction
of functions from topology because if $W \subset V \subset U$
and $s$ is a section of $\mathcal{F}$ over $U$ then
$s|_W = (s|_V)|_W$ by the property of the restriction maps
expressed in the definition above.

\medskip\noindent
Another notation that is often used is to indicate sections
over an open $U$ by the symbol $\Gamma(U, -)$ or by
$H^0(U, -)$. In other words, the following equalities
are tautological
$$
\Gamma(U, \mathcal{F}) = \mathcal{F}(U) = H^0(U, \mathcal{F}).
$$
In this chapter we will not use this notation, but in others
we will.

\begin{definition}
\label{definition-constant-presheaf}
Let $X$ be a topological space. Let $A$ be a set.
The {\it constant presheaf with value $A$} is the
presheaf that assigns the set $A$ to every open
$U \subset X$, and such that all restriction mappings
are $\text{id}_A$.
\end{definition}
